% Evaluation: the prospective confirmatory mapping protocol.
The empirical claim in this paper is governed by an execution-frozen protocol
that was created after the initial 32-case result but before any confirmatory
system output was examined. The predecessor result was allowed to motivate
replication; it was not allowed to change the second holdout, the evaluator,
the margins, or the pass criteria. The broader raw-text and expert-gold design
remains a separate, unexecuted study. Both freezes are identified in
Section~\ref{sec:availability}.

\subsection{Primary confirmatory hypothesis}

The frozen replication hypothesis has two non-compensatory conditions:

\begin{description}
    \item[False-merge superiority:] the coordinate-governed-minus-flat
        false-merge point estimate must be at most $-0.05$, and the upper
        bound of the paired 95\% bootstrap interval must be below zero.
    \item[False-split non-inferiority:] the
        coordinate-governed-minus-exact-conservative false-split point
        estimate must be at most $+0.03$, and the upper bound of its paired
        95\% bootstrap interval must be at most $+0.03$.
\end{description}

The holdout passes only if both conditions hold. It contains 32 cases, has
zero overlap with the initial 32-case set, and was bound by exact content
digest before any confirmatory outcome was accessed.

\subsection{Compared mapping rules}

The confirmatory primary comparison uses three deterministic rules over the
same already-structured mapping cases:

\begin{enumerate}
    \item \textbf{Typed mapping calculus}: evaluates the scientific
        coordinates the case requires and can retain an obstruction rather
        than force a merge.
    \item \textbf{Flat predicate canonicalization}: treats compatible surface
        predicates as sufficient merge evidence without the coordinate
        distinctions. This is the predeclared false-merge comparator.
    \item \textbf{Exact-coordinate conservative control}: requires exact
        coordinate compatibility and can abstain where it lacks such a match.
        This is the predeclared false-split and non-inferiority control.
\end{enumerate}

These are mapping-layer comparators, not attempts to reproduce a full external
knowledge-graph, retrieval-augmented or schema-fusion system. Stronger
schema-induction and provenance-contract baselines would be useful
post-outcome construct-validity pressure, but they cannot rewrite a
confirmatory status that was frozen beforehand.

\subsection{Secondary ablations}

The frozen mapping evaluator also reports descriptive paired ablations. The
covered attacks are:

\begin{itemize}
    \item \textbf{force compatibility without obstruction}: removes the
          explicit non-merge terminal and chooses a compatible mapping even
          where the coordinates prohibit it;
    \item \textbf{remove modality/polarity/attribution/discourse}: removes the
          coordinate group carrying the main confirmatory discrimination;
    \item \textbf{remove referent}, \textbf{construct},
          \textbf{measurement}, and \textbf{temporal context}: retained as
          zero-effect controls on this particular holdout.
\end{itemize}

The ablations are secondary and descriptive. The protocol explicitly forbids
turning their intervals into new confirmatory significance claims, and forbids
reading zero effect on an under-supported stratum as evidence that a
coordinate is unnecessary.

\subsection{Metrics and statistics}

The confirmatory decision uses case-level false-merge and false-split rates.
Wilson 95\% intervals describe individual rates. Paired differences use a
10,000-resample percentile bootstrap with frozen seed 20260817. The mapping
rules are deterministic, so the protocol uses one execution per case rather
than artificial stochastic repeats. Results are reported pooled and
descriptively by case family and source stratum. The initial 32 cases are
never pooled into the confirmatory verdict.

The primary decision rule is therefore fixed entirely by the two paired
comparisons above. No post-outcome margin change, exclusion, reweighting, or
replacement baseline can change whether the confirmatory run passed.

\subsection{Custody and fail-closed rules}

Before any confirmatory output was produced, the execution manifest bound the
portable confirmatory gold digest, zero overlap with the predecessor gold, the
exact MUSE, SciFact and SciSchema source revisions, the evaluator component
identities, the bootstrap seed and pass margins, and a rule that any digest
mismatch, source drift, evaluator drift or case overlap leaves the result
undetermined rather than silently re-evaluating a different study.

The gold is not an input to the evaluated mapping rule. The evaluator consumes
the candidate mapping result and the protected expected decision only after
the system output has been sealed.

\subsection{Reproducibility}

The result is reproduced by deterministic public-reference scripts over
immutable archives: the confirmatory analysis, publication figures and tables,
source registry, evaluator identities, and an independent replay receipt are
all retained, and two make targets regenerate the headline analysis and the
publication artifacts from committed inputs. Section~\ref{sec:availability}
names each of them.

This reproducibility claim does not extend to the unexecuted raw-text expert
atlas, to full external retrieval or schema baselines, to generated-portrait
recoverability, or to downstream answer quality. Each of those remains
undetermined until its own prospective run exists.

\endinput

Repository-level integrity is not scientific authority. The hashes, frozen
digests, deterministic replays and receipt chains this paper relies on establish
that a stated computation was the computation performed and can be performed
again. They establish nothing about whether its conclusion is correct, whether
the benchmark measures what it claims, or whether an independent party has
reviewed either. A perfectly replayed run of a flawed protocol replays the flaw
exactly. Where this paper's claims require external scientific adjudication,
that adjudication has not occurred and the integrity record does not substitute
for it.
